\section{Mechanism-structure-aware sensor fusion}

\todo[inline]{3 zdania o fuzji danych - uzyskanie pelnej informacji na podst. szczątkowych danych}

\subsection{The Extended Kalman Filter}

\todo[inline]{3 zdania o filtrze EKF - równanie głowne, oznaczenia}

\subsection{Constraints in multi-body systems}

\todo[inline]{Ograniczenia występujące w robotyce - wynikające z własciwosci konstrucyjnych itd}

\subsection{State constraints}

\todo[inline]{Orginalna metoda Dana Simona na 3 faze filtra EKF (fazę constraint correction) tj jak zadbać o to zeby stan spełniał dodatkowe równania}

\todo[inline]{Porównanie do metody N-R}

\subsection{Improvement %TODO}

\todo[inline]{Skoro podobne do N-R to dodajmy modyfikację w postacji alfa i k - poprawkę skalujemy x alfa (damped Newton's method) i zwiększamy liczbę iteracji na jeden obrót filtra}

